\labsection{9}{File-system namespace, paths, and links}{sec:lab09}
\begin{labproject}{Build a directory tree and investigate names, inode identity, and links}
\textbf{Suggested time:} 90--100 minutes \hfill
\textbf{Level:} introductory/intermediate\par
\textbf{Starter files:} \texttt{make\_demo\_tree.sh}, \texttt{directory\_walk.c}.\par
\textbf{Evidence to keep:} hand-drawn tree, pathname examples, inode/link table, and recursive-walk output.

\begin{labmode}
Use the real Linux/WSL file system. No graphics library is required. The goal is to understand what a directory tree represents, not to draw a decorative visualization.
\end{labmode}

\medskip\noindent\textbf{Learning outcomes.}
\begin{enumerate}
  \item distinguish absolute and relative pathnames;
  \item explain why a filename is not the same thing as the file's metadata/inode;
  \item distinguish a hard link from a symbolic link at an introductory level;
  \item explain how recursive directory traversal follows the namespace.
\end{enumerate}

\medskip\noindent\textbf{Part A --- Create and draw a known hierarchy.}
\begin{lstlisting}[style=osbash]
cd experiments/Lab09
bash make_demo_tree.sh
find demo_tree -maxdepth 4 -print
\end{lstlisting}
Draw the tree by hand. Label directories, regular files, the hard link, and the symbolic link. Your diagram should reflect the names reported by \texttt{find}.

\medskip\noindent\textbf{Part B --- Navigate using different path forms.}
Move into \texttt{demo\_tree/home/student}. Display \texttt{notes.txt} using: (1) a relative path, (2) an absolute path, and (3) a relative path containing \texttt{..}. Run \texttt{realpath notes.txt} and explain what it resolves.

\medskip\noindent\textbf{Part C --- Compare names with inode identity.}
\begin{lstlisting}[style=osbash]
ls -li notes.txt notes-hardlink.txt code-link.c
stat notes.txt
stat notes-hardlink.txt
readlink code-link.c
\end{lstlisting}
Record inode numbers and link counts. Explain why the hard-linked names refer to the same inode while a symbolic link stores a pathname to another object.

\begin{keyidea}
A directory gives names to objects in a namespace. The inode/metadata identifies properties of the underlying file object; storage blocks hold the file's data. These are related layers, not the same thing.
\end{keyidea}

\medskip\noindent\textbf{Part D --- Walk the tree recursively.}
\begin{lstlisting}[style=osbash]
gcc -Wall -Wextra directory_walk.c -o directory_walk
./directory_walk demo_tree
\end{lstlisting}
Find the recursive function call in the source. Explain the stopping rule and why the program deliberately does not recurse through symbolic links.

\medskip\noindent\textbf{Part E --- Controlled namespace change.}
Rename \texttt{notes.txt} while the hard link still exists. Inspect both names again. Then remove one hard link and inspect the remaining link count. Record the observations before deleting the demo tree.

\begin{verification}
Submit the tree, three pathname examples, the inode/link comparison, and a short explanation of why following symbolic links recursively without a policy could create cycles.
\end{verification}
\end{labproject}
